\newpage
\section{\tkzClass{square} 类}
变量 \tkzVar{square}{S} 持有一个表，用于存放正方形。它是可选的，你可以自由选择变量名。不过，为清晰和一致起见，推荐使用 \code{S}。如果你使用自定义变量（例如 Squares），则必须手动初始化它。如果使用了 \code{S}，\code{init\_elements()} 函数会重新初始化该表。


\subsection{创建正方形}
\label{sub:creating_a_square}

\tkzClass{square} 类根据两个相邻顶点构造一个正方形。点的顺序决定了方向。


\medskip
结果存放在 \tkzVar{square}{S} 中。

\begin{mybox}
\begin{verbatim}
S.ABCD = square:new(z.A, z.B)
\end{verbatim}
\end{mybox}

\textbf{简写形式。}

简写形式 \code{square(z.A, z.B)} 与之等价：

\begin{mybox}
\begin{verbatim}
S.ABCD = square(z.A, z.B)
\end{verbatim}
\end{mybox}

\subsection{正方形的属性}
\label{sub:square_attributes}
各点按正向（逆时针）给出。不做检验：四个点必须构成一个正方形。
\begin{mybox}
创建 | S.AB = square(z.A,z.B,z.C,z.D)|
\end{mybox}


  \bgroup
  \catcode`_=12
  \small
  \captionof{table}{正方形的属性。}\label{square:attributes}
  \begin{tabular}{lll}
  \toprule
  \textbf{属性}        & \textbf{用法}  &  \\
  \midrule
  \tkzAttr{square}{pa}         & |z.A = S.AB.pa|       &  \\
  \tkzAttr{square}{pb}         & |z.B = S.AB.pb|       &  \\
  \tkzAttr{square}{pc}         & |z.C = S.AB.pc|       &  \\
  \tkzAttr{square}{pd}         & |z.D = S.AB.pd|       &  \\
  \tkzAttr{square}{type}       & |S.AB.type= 'square'| &  \\
  \tkzAttr{square}{side}       & |s = S.AB.side|       & s = 边长   \\
  \tkzAttr{square}{center}     & |z.I = S.AB.center|   & 正方形的中心 \\
  \tkzAttr{square}{circumradius}   & |S.AB.circumradius|       & 外接圆半径 \\
  \tkzAttr{square}{inradius}   & |S.AB.inradius|      & 内切圆半径（边心距）  \\
  \tkzAttr{square}{apothem}    & |S.AB.apothem|       & 与内切圆半径相同  \\
  \tkzAttr{square}{apothem\_foot}       & |S.AB.apothem_foot|   & 中心在某条边上的投影 \\
  \tkzAttr{square}{diagonal}   & |S.AB.diagonal|      & 对角线长度  \\
  \tkzAttr{square}{area}       & |S.AB.area|          & 面积  \\
  \tkzAttr{square}{perimeter}  & |S.AB.perimeter|     & 周长  \\
  \tkzAttr{square}{ab}         & |S.AB.ab|             &  过两个顶点的直线   \\
  \tkzAttr{square}{ac}         & |S.AB.ac|             &  同上。 \\
  \tkzAttr{square}{da}         & |S.AB.da|             &  同上。 \\
  \tkzAttr{square}{bc}         & |S.AB.bc|             &  同上。 \\
  \tkzAttr{square}{bd}         & |S.AB.bd|             &  同上。 \\
  \tkzAttr{square}{cd}         & |S.AB.cd|             &  同上。 \\
  \bottomrule %
  \end{tabular}
  \egroup





\subsubsection{示例：正方形的属性}
\label{ssub:example_square_attributes}

\begin{minipage}{.5\textwidth}
\begin{verbatim}
\directlua{
 init_elements()
 z.A = point(0, 0)
 z.B = point(4, 0)
 z.C = point(4, 4)
 z.D = point(0, 4)
 S.new = square(z.A, z.B, z.C, z.D)
 z.I = S.new.center
 z.H = S.new.proj}
\begin{tikzpicture}
\tkzGetNodes
\tkzDrawCircles[orange](I,A I,H)
\tkzDrawPolygon(A,B,C,D)
\tkzDrawPoints(A,B,C,D,H,I)
\tkzLabelPoints(A,B,H,I)
\tkzLabelPoints[above](C,D)
\tkzDrawSegments(I,B I,H)
\tkzLabelSegment[sloped](I,B){%
  \pmpn{\tkzUseLua{S.new.circumradius}}}
\tkzLabelSegment[sloped](I,H){%
  \pmpn{\tkzUseLua{S.new.inradius}}}
\tkzLabelSegment[sloped](D,C){%
  \pmpn{\tkzUseLua{S.new.side}}}
\end{tikzpicture}
\end{verbatim}
\end{minipage}
\begin{minipage}{.5\textwidth}
\directlua{
 init_elements()
 z.A = point(0, 0)
 z.B = point(4, 0)
 z.C = point(4, 4)
 z.D = point(0, 4)
 S.new = square(z.A, z.B, z.C, z.D)
 z.I = S.new.center
 z.H = S.new.proj}

\begin{center}
  \begin{tikzpicture}
  \tkzGetNodes
  \tkzDrawCircles[orange](I,A I,H)
  \tkzDrawPolygon(A,B,C,D)
  \tkzDrawPoints(A,B,C,D,H,I)
  \tkzLabelPoints(A,B,H,I)
  \tkzLabelPoints[above](C,D)
  \tkzDrawSegments(I,B I,H)
  \tkzLabelSegment[sloped](I,B){\pmpn{\tkzUseLua{S.new.circumradius}}}
  \tkzLabelSegment[sloped](I,H){\pmpn{\tkzUseLua{S.new.inradius}}}
  \tkzLabelSegment[sloped](D,C){\pmpn{\tkzUseLua{S.new.side}}}
  \end{tikzpicture}
\end{center}
\end{minipage}


\newpage



\subsection{正方形的方法与函数}
\label{sub:square_methods}


  \bgroup
  \catcode`_=12
  \small
  \captionof{table}{正方形的方法}\label{square:methods}
  \begin{tabular}{lll}
  \toprule
  \textbf{方法} & \texttt{参考}&    \\
  \midrule

  \tkzMeth{square}{new(za,zb,zc,zd)} &[\ref{ssub:method_square_rotation}]\\

  \tkzFct{square}{square.by\_rotation (zi,za)} &[\ref{ssub:method_square_rotation}]\\

  \tkzFct{square}{square.from\_side(za,zb,"swap")} & \ref{ssub:square_with_side_method} \\
  \tkzMeth{square}{get(<i>)} & [\ref{ssub:method_square_get}]\\
  \bottomrule %
  \end{tabular}
  \egroup




\subsubsection{方法 \tkzMeth{square}{get(<i>)}}
\label{ssub:method_square_get}

该方法返回定义正方形对象的各个点。

\begin{itemize}
  \item \code{S.ABCD:get()} 按顺序返回四个顶点：\code{pa}、\code{pb}、\code{pc}、\code{pd}；
  \item \code{S.ABCD:get(i)}（\code{i} 取 1 到 4）只返回第 \code{i} 个顶点；例如 \code{S.ABCD:get(2)} 等价于 \code{S.ABCD.pb}。
\end{itemize}

\begin{tkzexample}[code only]
\directlua{
  init_elements()
  z.A = point(0, 0)
  z.B = point(2, 0)
  S.ABCD = square.from_side(z.A, z.B)
  z.E, z.F, z.G, z.H = S.ABCD:get()
  z.K = S.ABCD:get(3)
}
\end{tkzexample}

\subsubsection{函数 \tkzFct{square}{square.by\_rotation(pt,pt)}}
\label{ssub:method_square_rotation}
$I$ 为正方形中心；$A$ 为第一个顶点

\vspace{1em}
\begin{tkzexample}[latex=.5\textwidth]
\directlua{
 z.A  = point(0, 0)
 z.I  = point(2, -1)
 S    = square.by_rotation(z.I, z.A)
 z.B  = S.pb
 z.C  = S.pc
 z.D  = S.pd
 z.I  = S.center}
\begin{tikzpicture}
\tkzGetNodes
\tkzDrawPolygon(A,B,C,D)
\tkzDrawPoints(A,B,C,D)
\tkzLabelPoints(B)
\tkzLabelPoints[above](A)
\tkzLabelPoints[right](C,D)
\tkzDrawPoints[red](I)
\end{tikzpicture}
\end{tkzexample}

\subsubsection{方法 \tkzFct{square}{square.from\_side(za,zb)}}
\label{ssub:square_with_side_method}
若带选项 \code{"swap"}，则正方形按逆时针方向定义。也可以由一条直线得到同样的结果 [\ref{ssub:method_line_square}]。

\begin{minipage}{.5\textwidth}
\begin{verbatim}
\directlua{
 init_elements()
 z.A = point(0, 0)
 z.B = point(2, 1)
 S.side = square.from_side(z.A, z.B)
 z.B = S.side.pb
 z.C = S.side.pc
 z.D = S.side.pd
 z.I = S.side.center}
\begin{tikzpicture}[scale = 2]
   \tkzGetNodes
   \tkzDrawPolygon(A,B,C,D)
   \tkzDrawPoints(A,B,C,D)
   \tkzLabelPoints(A,B)
   \tkzLabelPoints[above](C,D)
   \tkzDrawPoints[red](I)
\end{tikzpicture}
\end{verbatim}
\end{minipage}
\begin{minipage}{.5\textwidth}
\directlua{
 init_elements()
 z.A = point(0, 0)
 z.B = point(2, 1)
 S.side = square.from_side(z.A, z.B)
 z.B = S.side.pb
 z.C = S.side.pc
 z.D = S.side.pd
 z.I = S.side.center}
\begin{center}
  \begin{tikzpicture}[scale = 2]
  \tkzGetNodes
  \tkzDrawPolygon(A,B,C,D)
  \tkzDrawPoints(A,B,C,D)
  \tkzLabelPoints(A,B)
  \tkzLabelPoints[above](C,D)
  \tkzDrawPoints[red](I)
  \end{tikzpicture}
\end{center}

\end{minipage}
\endinput